% !TEX program = xelatex
\documentclass[12pt,a4paper]{report}

\usepackage{graphicx}
\usepackage{geometry}
\usepackage{float}
\usepackage{xcolor}
\usepackage{titlesec}
\usepackage[space]{grffile}
\usepackage{xepersian}

\geometry{top=3cm, bottom=3cm, left=2.5cm, right=3cm}

% تنظیم فاصله بعد از عنوان‌ها
\titlespacing*{\section}{0pt}{2ex plus 1ex minus .2ex}{1.5ex plus .2ex}
\titlespacing*{\subsection}{0pt}{1.8ex plus 1ex minus .2ex}{1ex plus .2ex}

% فونت فارسی ۱۴ و انگلیسی ۱۲
\settextfont[
Path = C:/USERS/SAZGAR/APPDATA/LOCAL/MICROSOFT/WINDOWS/FONTS/,
BoldFont = IRNAZANIN.TTF,
Scale = 1.166
]{IRNAZANIN.TTF}
\setlatintextfont[Scale=1]{Times New Roman}

\linespread{1.5}

\begin{document}
	\sloppy
	
	% ----------------------------------------------------------------
	% فصل ۳: تهیه دادگان (بدون شماره خودکار)
	% ----------------------------------------------------------------
	
	\chapter*{\bfseries فصل ۳: تهیه دادگان}
	\addcontentsline{toc}{chapter}{فصل ۳: تهیه دادگان}
	
	\vspace{-2.5em}
	\noindent\rule{\textwidth}{0.5pt}
	\vspace{0.5em}
	
	% ----------------------------------------------------------------
	% بخش مقدمه
	% ----------------------------------------------------------------
	\section{مقدمه}
	
	توسعه مدل‌های زبانی بزرگ \lr{(LLMs)} در سال‌های اخیر تحولی بنیادین در پردازش زبان طبیعی ایجاد کرده است، اما کارایی این مدل‌ها به شدت وابسته به کیفیت و تنوع داده‌هایی است که با آن‌ها آموزش دیده یا ارزیابی می‌شوند. در حوزه تخصصی حقوق، این وابستگی اهمیتی دوچندان می‌یابد؛ چرا که متون حقوقی نه تنها از نظر واژگانی پیچیده هستند، بلکه ساختاری منطقی، سلسله‌مراتب‌دار و به‌شدت وابسته به بافت \lr{(Context-Dependent)} دارند. در اکوسیستم حقوقی ایران، چالش‌های زبانی زبان فارسی با پیچیدگی‌های ذاتی نظام حقوقی درهم آمیخته است و فقدان یک پیکره \lr{(Corpus)} جامع، ساختاریافته و استاندارد، همواره گلوگاه اصلی پژوهش‌های هوش مصنوعی در این حوزه بوده است.
	
	بسیاری از مدل‌های زبانی عمومی که بر روی داده‌های وب فارسی آموزش دیده‌اند، اگرچه در تولید متن روان موفق عمل می‌کنند، اما در مواجهه با «استدلال حقوقی» و «استناد دقیق به مواد قانونی» دچار توهم \lr{(Hallucination)} یا خطاهای فاحش می‌شوند. این ضعف ناشی از آن است که داده‌های عمومی فاقد عمق معنایی لازم برای درک مفاهیم حقوقی هستند. برای غلبه بر این چالش و پر کردن شکاف «دانش تخصصی»، در این پژوهش تمرکز ویژه‌ای بر گردآوری یک مجموعه داده چندوجهی و غنی شده است.
	
	رویکرد ما در تهیه دادگان، فراتر از جمع‌آوری صرف متون خام بوده است. ما با درک این نکته که دانش حقوقی تنها در «متن قانون» خلاصه نمی‌شود، بلکه در «رویه قضایی» (نحوه تفسیر قانون توسط قضات) و «تعاملات واقعی» (سوالات شهروندان و پاسخ وکلا) جاری است، استراتژی جمع‌آوری داده‌ها را بر سه رکن اصلی بنا نهادیم:
	
	\textbf{۱. قوانین و مقررات رسمی (منابع آزمون وکالت)}
	
	این بخش، زیربنای اصلی دانش مدل را تشکیل می‌دهد. برای اطمینان از جامعیت و اعتبار داده‌ها، دقیقاً از همان مجموعه قوانینی استفاده کردیم که به عنوان منابع رسمی آزمون وکالت معرفی شده‌اند. این انتخاب هوشمندانه باعث می‌شود مدل بر روی مهم‌ترین و کاربردی‌ترین قوانین کشور (مانند قانون مدنی، تجارت، مجازات اسلامی و...) متمرکز شود. این متون به مدل می‌آموزند که «قانون دقیقاً چه می‌گوید» و مواد قانونی چه ساختار و کلماتی دارند.
	
	\textbf{۲. رویه قضایی و آرای دادگاه‌ها}
	
	دانستن متن قانون کافی نیست؛ مدل باید یاد بگیرد که قاضی‌ها چگونه از این قوانین برای حل پرونده‌های واقعی استفاده می‌کنند. برای این منظور، از دو نوع داده استفاده کردیم:
	
	\textbf{آرای وحدت رویه:} گاهی اوقات دادگاه‌های مختلف در موارد مشابه، برداشت‌های متفاوتی از یک قانون دارند. در این شرایط، دیوان عالی کشور یا دیوان عدالت اداری رایی صادر می‌کند که برای همه دادگاه‌ها لازم‌الاجراست و حکم قانون را دارد. ما این آرا را جمع‌آوری کردیم تا مدل با تفسیرهای قطعی و نهایی قانون آشنا شود.
	
	\textbf{دادنامه‌های قضایی:} این‌ها همان آرای معمولی دادگاه‌ها در پرونده‌های روزمره هستند. ما تعداد زیادی از این آرا را جمع‌آوری کردیم تا مدل با استدلال‌های قضایی، نحوه نگارش رای توسط قاضی و چگونگی تطبیق قانون با وقایع پرونده آشنا شود.
	
	\textbf{۳. پرسش و پاسخ‌های حقوقی (زبان طبیعی مردم)}
	
	مردم عادی معمولاً با زبان تخصصی حقوقی صحبت نمی‌کنند؛ آن‌ها مشکلاتشان را با زبان ساده و روزمره بیان می‌کنند، برای اینکه مدل بتواند این زبان را بفهمد و به درستی راهنمایی کند، از مجموعه‌ای از سوالات واقعی مردم و پاسخ‌های وکلا استفاده کردیم. این بخش به مدل کمک می‌کند تا مثل یک مترجم عمل کند: مشکل را به زبان ساده بشنود و راه‌حل را با استناد به قانون ارائه دهد.
	
	\section{منابع و مجموعه داده}
	
	در راستای پیاده‌سازی رکن اول، تمرکز بر قوانینی قرار گرفت که مبنای اصلی آزمون وکالت در ایران هستند. این انتخاب از آن جهت حائز اهمیت است که منابع آزمون وکالت، کاربردی‌ترین قوانین کشور را شامل می‌شوند. متن کامل این قوانین از پایگاه اطلاع‌رسانی حقوقی «اختبار» استخراج گردید؛ این پایگاه به دلیل دسترسی‌پذیری عمومی و به‌روزرسانی مداوم اصلاحیه قوانین، بستری قابل اتکا برای دریافت متن خام مواد قانونی محسوب می‌شود. از حوزه‌های حقوقی مورد استفاده می‌توان به حقوق مدنی، آیین دادرسی مدنی، حقوق تجارت، حقوق جزای عمومی، حقوق جزای اختصاصی، آیین دادرسی کیفری و حقوق اساسی اشاره کرد.
	
	برای پوشش رکن دوم و آموزش نحوه استدلال قضایی به مدل، نیاز به دسترسی به مخزنی غنی از آرای صادر شده توسط مراجع عالی و تالی قضایی بود. بدین منظور، از آرشیو دیجیتال وب‌سایت «داودآبادی» بهره‌برداری شد. این منبع ارزشمند، دسترسی به دو دسته حیاتی از داده‌ها را فراهم کرد: نخست، آرای وحدت رویه صادر شده از سوی دیوان عالی کشور و دیوان عدالت اداری که حکم قانون را داشته و فصل‌الخطاب تعارضات قضایی هستند؛ و دوم، مجموعه وسیعی از دادنامه‌های انشایی صادر شده از شعب مختلف دادگاه‌ها که شیوه نگارش رأی، استناد به مواد و تطبیق حکم با موضوع را به تصویر می‌کشند.
	
	در فرایند جمع‌آوری آرای وحدت رویه، با مشورت یک وکیل متخصص، آرای دیوان عدالت اداری از سال ۱۳۹۰ به بعد و آرای دیوان عالی کشور از سال ۱۳۶۶ به بعد برای استخراج انتخاب گردید. این بازه زمانی به دلیل تغییرات ساختاری قوانین و اهمیت آرای دوره‌های اخیر تعیین شد. در نهایت، حدود ۲۷٬۹۳۹ دادنامه قضایی استخراج و پردازش شد که این حجم از داده، پوشش جامعی از رویه‌های قضایی در حوزه‌های مختلف حقوقی را فراهم آورد.
	
	در نهایت، برای تحقق رکن سوم و پر کردن شکاف میان زبان تخصصی و زبان عامیانه، از مجموعه داده جامع «پرسش و پاسخ حقوقی ایران» \lr{(Iranian Legal Question Answering Dataset)} موجود در مخزن هاگینگ‌فیس استفاده شد. مرجع اصلی گردآوری این داده‌ها، وب‌سایت مشاوره حقوقی «دادراه» است که تعاملات واقعی میان مردم و وکلا را میزبانی می‌کند. این مجموعه داده عظیم شامل بیش از ۶۰۰ هزار رکورد پرسش است که توسط گویشوران عادی زبان فارسی با ادبیاتی غیررسمی مطرح شده و در مقابل، بیش از ۲ میلیون پاسخ تشریحی از سوی وکلای دادگستری برای آن‌ها ثبت شده است. ساختار غنی این دادگان که شامل فیلدهای کلیدی نظیر «عنوان سوال» \lr{(Question Title)}، «متن کامل بدنه سوال» (\lr{Question Body})، «پاسخ‌ها» \lr{(Answers)} و «برچسب‌های موضوعی» \lr{(Tags)} است، به مدل اجازه می‌دهد تا الگوی معنایی تبدیل دغدغه‌های روزمره شهروندان به مفاهیم انتزاعی حقوقی را فراگرفته و با تنوع زبانی گسترده‌ای از سناریوهای حقوقی آشنا شود.
	
	\section{فرآیند گردآوری و پیش‌پردازش دادگان}
	
	پس از شناسایی منابع معتبر، گام عملیاتی جهت تبدیل این منابع به یک پایگاه داده ساختاریافته آغاز گردید. این فرایند در دو فاز اصلی «استخراج داده» \lr{(Data Extraction)} و «پیش‌پردازش و غنی‌سازی» (\lr{Preprocessing \& Enrichment}) انجام شد که جزئیات فنی هر مرحله در ادامه تشریح شده است.
	
	\subsection{مرحله اول: استخراج داده‌ها از طریق خزش وب}
	
	با توجه به عدم دسترسی به API رسمی برای قوانین و رویه‌های قضایی، از تکنیک خزش وب برای جمع‌آوری داده‌ها استفاده شد. بدین منظور، پردازه‌نویسی‌های اختصاصی برای هر کدام از منابع هدف (سایت‌های اختبار و داودآبادی) توسعه یافت تا ساختار سلسله‌مراتبی صفحات وب پیمایش و محتوای هدف استخراج گردد. 
	
	الگوریتم خزش به گونه‌ای طراحی شد که نه تنها «متن اصلی» قوانین (منابع آزمون وکالت)، بلکه کلیه ملحقات ضروری مانند آیین‌نامه‌ها، دستورالعمل‌های اجرایی و تبصره‌های مرتبط را نیز شناسایی و ذخیره کند. همچنین در حوزه رویه قضایی، تمام آرای وحدت رویه (دیوان عالی و دیوان عدالت) و دادنامه‌های منتخب دادگاه‌ها به همراه فراداده‌های موجود، استخراج و در مخزن اولیه بارگذاری شدند.
	
	آرای وحدت رویه دیوان عدالت اداری سال ۱۳۹۰ به بعد و آرای وحدت رویه دیوان عالی سال ۱۳۶۶ به بعد، پس از مشورت با یک وکیل متخصص، به‌عنوان ملاک استخراج انتخاب شدند. این انتخاب در راستای تمرکز بر رویه‌های به‌روز و کاربردی صورت گرفت. در مجموع، حدود ۲۷٬۹۳۹ دادنامه استخراج گردید.
	
	\subsection{مرحله دوم: پیش‌پردازش و استانداردسازی متن}
	
	داده‌های استخراج‌شده از وب معمولاً دارای نویز و ناهمگونی‌های متعددی بودند. برای آماده‌سازی این متون جهت آموزش مدل، یک خط لوله پیش‌پردازش متن با تاکید بر ویژگی‌های زبان فارسی و الزامات متون حقوقی اجرا شد:
	
	\textbf{اصلاح نویسه‌ها:} تمامی کاراکترهای عربی (مانند «ي» و «ك») به معادل فارسی جایگزین گردیدند. همچنین اعداد، چه در قالب لاتین و چه عربی، به فرمت استاندارد فارسی یکدست شدند.
	
	\textbf{مدیریت نیم‌فاصله:} برخلاف بسیاری از پردازش‌های متن، به دلیل اهمیت ترکیبات واژگانی در حقوق، تمامی نیم‌فاصله‌ها با دقت ترمیم و حفظ گردید تا یکپارچگی کلمات مرکب رعایت شود.
	
	\subsection{مرحله سوم: استخراج ساختار و فراداده}
	
	ارزش اساسی یک دادگان حقوقی فقط در متن خام نیست، بلکه در «ساختار» و «فراداده‌» آن نهفته است. بدین منظور، غنی‌سازی داده‌ها در دو سطح انجام شد:
	
	\textbf{الف) ساختاردهی قوانین:}
	
	هر قانون دارای یک سلسله‌مراتب درختی پیچیده است. برای حفظ این ساختار، الگوریتم تجزیه‌گر \lr{(Parser)} توسعه‌یافته، متن قوانین را پیمایش کرده و اجزای ساختاری آن شامل «جلد»، «کتاب»، «بخش»، «باب»، «فصل»، «مبحث» و «فقره» را شناسایی کرد. این اطلاعات در ستون‌های مجزا ذخیره شدند تا امکان بازیابی دقیق هر ماده قانونی بر اساس جایگاه آن در ساختار قانون فراهم شود. علاوه بر این، نام کامل قانون و کلیدواژه‌های اصلی مرتبط با آن نیز به هر سطر \lr{(Record)} اضافه شد.
	
	\textbf{ب) غنی‌سازی آرای قضایی:}
	
	برای آرای وحدت رویه و دادنامه‌ها، صرفِ متن رای کافی نبود. با استفاده از تکنیک‌های تطبیق الگو و پردازش متن، اطلاعات کلیدی شامل «شماره دادنامه»، «تاریخ صدور رای» و «شعبه صادرکننده» از متن استخراج شد. همچنین با تحلیل محتوای رای و تطبیق آن با لیست قوانین پرکاربرد، «حوزه تخصصی» (مانند مدنی، کیفری، خانواده و...) و «کلمات کلیدی» مرتبط شناسایی و به عنوان متادیتا به هر رای الصاق گردید. این برچسب‌گذاری معنایی، قابلیت جستجو و آموزش تخصصی مدل را به طرز چشمگیری افزایش می‌دهد.
	
	\section{جمع‌بندی فصل}
	
	در این فصل، مسیر تبدیل منابع پراکنده حقوقی به یک پیکره تمیز و ماشین‌خوان تشریح شد. ما با عبور از سه مرحله کلیدی «شناسایی منابع معتبر»، «استخراج هوشمند داده‌ها» و «پالایش و ساختاردهی دقیق»، موفق به ایجاد مجموعه‌ای شدیم که نه تنها حجم قابل‌توجهی از دانش حقوقی ایران را در بر می‌گیرد، بلکه با حفظ ساختارهای سلسله‌مراتبی و افزودن فراداده‌های غنی، آمادگی لازم برای آموزش مدل‌های زبانی پیچیده را داراست. این دادگان، زیربنای اصلی آزمایش‌ها و ارزیابی‌هایی است که در فصول آینده به آن‌ها خواهیم پرداخت.
	
\end{document}
